\newpage
\phantomsection
\label{prf:limit-difference}

\begin{remark*}[Return]
\hyperref[thm:limit-difference]{Return to Theorem}
\end{remark*}

\begin{theorem*}[Difference Rule for Function Limits]
Let $f, g : A \to \mathbb{R}$, let $c$ be a cluster point of $A$. If
$\lim_{x\to c}f(x)=L$ and $\lim_{x\to c}g(x)=M$, then
$\lim_{x\to c}(f-g)(x) = L-M$.
\end{theorem*}

\begin{proof}
\textbf{Professional Standard Proof.}~\\
Let $(x_n) \subseteq A \setminus \{c\}$ with $x_n \to c$.
By the Sequential Criterion, $f(x_n) \to L$ and $g(x_n) \to M$.
By the difference rule for sequence limits, $(f-g)(x_n) = f(x_n)-g(x_n) \to L-M$.
Since $(x_n)$ was arbitrary, the Sequential Criterion gives
$\lim_{x\to c}(f-g)(x) = L-M$.
\end{proof}

\begin{proof}
\textbf{Detailed Learning Proof.}~\\
\textbf{Step 1. Set up a test sequence.}
Let $(x_n) \subseteq A \setminus \{c\}$ be any sequence with $x_n \to c$.

\textbf{Step 2. Apply the Sequential Criterion to each component.}
Since $\lim_{x\to c}f(x)=L$, the Sequential Criterion gives $f(x_n) \to L$.
Since $\lim_{x\to c}g(x)=M$, the Sequential Criterion gives $g(x_n) \to M$.

\textbf{Step 3. Apply the difference rule for sequences.}
The algebra of sequence limits gives
\[
  (f-g)(x_n) = f(x_n) - g(x_n) \to L - M.
\]

\textbf{Step 4. Lift back by the Sequential Criterion.}
Since every such $(x_n)$ satisfies $(f-g)(x_n) \to L-M$, the Sequential
Criterion yields $\lim_{x\to c}(f-g)(x) = L-M$.
\end{proof}

\begin{remark*}[Proof structure]
Bridge argument identical to the sum rule, with the sequence difference rule
in place of the sequence sum rule. The difference rule reduces to $f + (-1)g$.
\end{remark*}

\begin{remark*}[Dependencies]
\begin{itemize}
  \item \hyperref[prop:sequential-criterion-limits]{Sequential Criterion for Function Limits}
  \item \hyperref[def:limit-function]{Limit of a Function}
\end{itemize}
\end{remark*}
